\phantomsection
\markboth{RÉSUMÉ}{}
\thispagestyle{frontmatterstyle}

\FrontTitleSummary{Résumé}

\begin{FrontTextBlock}
Ce Projet de Fin d’Études porte sur la conception d’une solution data\mbox{-}driven pour le suivi et l’analyse des opérations de merchandising terrain du projet Unilever chez Smollan Morocco. Le travail part d’un processus initial fortement dépendant des exports \textit{Data Dump} et \textit{Coverage}, des traitements Excel et de rapports préparés séparément.

La solution proposée centralise les données dans une base MySQL à travers un pipeline Data Engineering, puis les exploite via des tableaux de bord Power BI, une application mobile de consultation côté client et un portail back-office interne. Un moteur de validation métier permet également d’identifier des situations nécessitant une revue, notamment à partir des contrôles OSA et GPS.

Les résultats obtenus montrent une amélioration de la traçabilité, une réduction du temps de préparation des données et une exploitation plus structurée des informations terrain. Le projet constitue ainsi une base fonctionnelle pour renforcer le suivi opérationnel et l’analyse des données merchandising dans un contexte professionnel réel.

\vspace{0.8cm}

\noindent\textbf{Mots-clés :} Data Engineering, ETL, MySQL, merchandising terrain, OSA, GPS, Power BI, application mobile, back-office, validation métier.
\end{FrontTextBlock}